We presented a real-hardware ROS2 harness for decoupled LLM/VLM
manipulation planning, built on the premise that a model's broad,
non-robot-specific pretraining is usable for physical reasoning through
code, without collecting task-specific action data, provided a
deterministic layer grounds the plan's numeric geometry in sensor
measurement rather than trusting the model's own arithmetic. Deploying
that harness on a real UR5cb manipulator turned this premise into an
observation: over six development sessions we found six recurring
geometry-grounding bugs, invisible to symbolic and scripted simulation by
construction, and closed five of them with the harness's code-grounding
layer. A 120-trial real-hardware evaluation, using a single unified model
for both grounding modalities, shows the sixth, a multi-object placement
collision, as the dominant residual failure, occurring in both the LLM
and VLM pipelines regardless of which one grounded the scene. Closing
that bug's dynamic, trajectory-level half as a new harness skill, along
with a depth-verified treatment of grasp width, are the most direct next
steps; more broadly, we argue that a geometry-grounding audit belongs
alongside task-complexity benchmarking for any LLM/VLM harness headed to
physical hardware.
